% ATTEMPT_STATUS: unresolved
% ATTEMPT_ORIGIN: new research, 2026-10-01
% TOP_PROBLEM: {"id": 1386, "title": "Danzer's Problem", "category_id": 6, "category": {"id": 6, "name": "geometry", "display_name": "Geometry", "description": "Euclidean and non-Euclidean geometry, geometric structures.", "slug": "geometry", "order_index": 6, "created_at": "2026-07-31T15:26:25.671Z"}, "status": "open"}
% FIRST_POSTED: 2026-10-01
% Imported from: unsolvedmath_additional_500.tex; UnsolvedMath ID 1386
% !TeX program = lualatex
% Compile twice: lualatex 1386.tex
% Self-contained: no external bibliography, images, or \input files.
% Numeric section labels and visible numbers are original UnsolvedMath IDs.
\documentclass[11pt,a4paper]{article}
\usepackage[margin=24mm,headheight=14pt]{geometry}
\usepackage{fontspec}
\setmainfont{Latin Modern Roman}
\setsansfont{Latin Modern Sans}
\setmonofont{Latin Modern Mono}
\usepackage{amsmath,amssymb,mathtools}
\usepackage{unicode-math}
\setmathfont{Latin Modern Math}
\usepackage{microtype}
\usepackage{enumitem,xurl,needspace}
\usepackage[dvipsnames]{xcolor}
\usepackage{hyperref}
\hypersetup{unicode=true,colorlinks=true,linkcolor=MidnightBlue,urlcolor=MidnightBlue,
 pdftitle={UnsolvedMath Problem 1386},pdfauthor={Source-annotated compilation},bookmarksnumbered=true}
\usepackage{bookmark,tocloft,titlesec,fancyhdr}
\setlength{\cftsecnumwidth}{4.2em}
\cftsetpnumwidth{2.5em}
\cftsetrmarg{4em}
\titleformat{\section}{\Large\bfseries\raggedright}{\thesection}{0.7em}{}
\pagestyle{fancy}\fancyhf{}
\fancyhead[L]{\small\sffamily UnsolvedMath research attempt}
\fancyhead[R]{\small Original catalogue IDs}
\fancyfoot[C]{\thepage}
\setlength{\parindent}{0pt}
\setlength{\parskip}{5pt plus 1pt}
\setlength{\emergencystretch}{3em}
\setcounter{tocdepth}{1}\setcounter{secnumdepth}{1}
\setlist[itemize]{leftmargin=3em,itemsep=4pt,topsep=4pt,parsep=0pt}
\allowdisplaybreaks
\newcommand{\N}{\mathbb{N}}
\newcommand{\Z}{\mathbb{Z}}
\newcommand{\Q}{\mathbb{Q}}
\newcommand{\R}{\mathbb{R}}
\newcommand{\C}{\mathbb{C}}
\newcommand{\F}{\mathbb{F}}
\newcommand{\A}{\mathbb{A}}
\newcommand{\PP}{\mathbb{P}}
\newcommand{\E}{\mathbb{E}}
\newcommand{\eps}{\varepsilon}
\newcommand{\dd}{\,\mathrm d}
\newcommand{\sref}[2]{\hyperlink{source:#1:#2}{[S#2]}}
\newif\ifshowcatalogue
\showcataloguetrue
\newcommand{\cataloguescope}[1]{\ifshowcatalogue
 \paragraph{Statement recorded in the supplied source.}
 #1\fi}
\begin{document}
% UnsolvedMath ID 1386; source catalogue code GEOM-027

\setcounter{section}{1385}

\section{Bounded-Density Danzer Sets in the Plane}\label{1386}

{\small\sffamily UnsolvedMath ID 1386 · Geometry}

\subsection{Definitions and mathematical statement}

\paragraph{Shared notation.}Unless a section says otherwise, $\N=\{1,2,\ldots\}$,
$\Z$ is the ring of integers, $\Q,\R,\C$ are the rational, real, and complex
fields, $[N]=\{1,\ldots,\lfloor N\rfloor\}$, and $\log$ is the natural
logarithm. For real functions of a parameter tending to infinity,
$f\ll g$ and $f=O(g)$ mean $|f|\leq Cg$ eventually for some fixed $C$;
$f\gg g$ reverses this comparison; $f=o(g)$ means $f/g\to0$;
$f\sim g$ means $f/g\to1$; and $f\asymp g$ means both order bounds.
A subscript on $O$ or $\ll$ specifies allowed constant dependence.
The expression $n^{o(1)}$ in an upper bound means growth smaller than
$n^\epsilon$ for each fixed $\epsilon>0$. Local definitions govern any
exception, finite-field convention, or use of the same symbol for another
object. An asymptotic claim is not automatically an effective algorithm.



A planar Danzer set $D$ meets every convex body of area one. The bounded-growth-density condition is $\#(D\cap B(0,R))=O(R^2)$ as $R\to\infty$. A uniformly separated set has some $\eta>0$ with $\|x-y\|_2\geq\eta$ for distinct $x,y\in D$. These are different requirements: a positive separation implies an upper density bound, but the converse need not hold. \sref{1386}{1} \sref{1386}{2}

\paragraph{Statement recorded in the supplied source.}
Do Danzer sets of bounded density or bounded separation exist? \sref{1386}{1}

\paragraph{Scope and interpretation.}

The two standard density and separation versions are distinguished. A uniformly bounded number of points in every convex unit-area body is a stronger, different condition and is not imposed. \sref{1386}{1}

\subsection{Short English statement}

Can a set meet every unit-area convex body without becoming too dense, or even while keeping a fixed positive separation between its points?

\subsection{Research attempt}
\textbf{Status: UNRESOLVED.} We prove that every full-rank lattice fails the Danzer property, distinguish separation from growth by explicit constructions, and solve a restricted bounded-shape version with a grid. The obstruction is the uniform treatment of arbitrarily thin convex bodies.

\paragraph{Source check and strategy.}
Solomon--Weiss [E1], consulted 1 October 2026, studies Danzer sets and dense forests, rules out several natural construction classes, and gives a logarithmic-loss growth bound. Its deep constructions are not inputs to the elementary arguments below. We test a lattice construction, isolate exactly how it fails, and then see which additional restriction repairs it. Intersecting every large disk is a weaker requirement than intersecting every unit-area convex body.

\paragraph{Every lattice has empty unit-area rectangles.}
First consider $D=\mathbb Z^2$. The closed rectangle $[0,2]\times[1/4,3/4]$ has area one and misses $D$, since its vertical coordinates lie strictly between consecutive integers. More generally let $L=A\mathbb Z^2$ with $A$ invertible. There is a nonzero linear functional $\ell$ taking integer values on $L$, for instance a row of $A^{-1}$. The open strip $1/4<\ell(x)<3/4$ contains no lattice point and has positive Euclidean width. Inside it choose a closed rectangle of any smaller positive width and reciprocal length. It has area one and misses $L$.

This proof works at every lattice scale. Making a lattice denser makes the empty strip thinner, but its rectangle can become proportionally longer. It also rules out every subset of a lattice, and finite unions of cosets that lie in a common sufficiently fine lattice. It does not rule out arbitrary finite unions of incommensurable lattice translates; their projections onto a selected normal direction may no longer leave an empty strip.

\paragraph{Separation implies the stated growth bound.}
If distinct points of $D$ are at least $\eta>0$ apart, the open disks of radius $\eta/2$ around points in $D\cap B(0,R)$ are disjoint and lie in $B(0,R+\eta/2)$. Comparing areas gives
\[
\#(D\cap B(0,R))\le\left(1+\frac{2R}{\eta}\right)^2.
\]
Thus a uniformly separated Danzer set would solve the bounded-growth version too. This implication uses no assumption about convex-body intersections.

For the converse, take $D=\{(m,0),(m,m^{-2}):m\ge2\}$. Its growth in centered balls is $O(R)$, hence $O(R^2)$, and it is locally finite. The two points associated to $m$ have distance $m^{-2}\to0$, so no positive separation constant exists. This set is not proposed as Danzer; it is a precise counterexample to identifying the two geometric hypotheses.

\paragraph{A restricted theorem for bodies of bounded roundness.}
Fix $K\ge1$ and consider only area-one convex bodies $C$ for which there are a centre $c$ and radius $r>0$ satisfying
\[
B(c,r)\subset C\subset B(c,Kr).
\]
The outer containment gives $1\le\pi K^2r^2$, so $r\ge1/(K\sqrt\pi)$. Every point of the plane is within $s/\sqrt2$ of the square lattice $s\mathbb Z^2$: round each coordinate to a nearest multiple of $s$. Choose $s<\sqrt2/(K\sqrt\pi)$. Then a lattice point lies strictly inside $B(c,r)$, and therefore in $C$.

This produces a uniformly separated hitting set for the restricted class, with an explicit spacing independent of the body's location or orientation. Its number of points in a radius-$R$ ball is $O(K^2R^2)$ by the preceding packing argument. Convexity is not even needed once the two containments are given; the inner-ball condition does the work.

\paragraph{Why taking a union over shape parameters fails.}
It is tempting to use finer and finer grids for $K=1,2,\ldots$ and unite them. Each individual grid has the needed guarantee for its own class, but the infinite union need not be locally finite. For example the union of grids with spacings $2^{-j}$ has infinitely many points in every disk with nonempty interior. Its counting function is infinite there, so it fails the required growth bound in the strongest possible way. Taking just one grid returns to the empty-strip obstruction.

Deleting points from fine grids to control their density could remove the very point that meets a thin body. A valid multiscale construction would need both a quantitative hitting proof and a summable local counting estimate. Neither follows by taking the union of the separate bounded-roundness solutions. This is a quantifier issue: one set must work for every shape parameter at once.

\paragraph{Verification and exact remaining gap.}
All rectangles used to refute lattices are closed and stay strictly inside an empty strip, so there is no boundary-point escape. The growth counterexample is locally finite and satisfies the centered-ball convention of the statement. The restricted theorem tolerates arbitrary translations and rotations; it is the roundness parameter that is bounded.

The unresolved step is constructing a single $O(R^2)$-growth set hitting every unit-area convex body, or proving such a set impossible. Uniform separation is an additional target. The calculations here exclude a natural family and identify a precise multiscale obstruction, but do not settle either full version. No new numerical search or literature-level novelty is claimed.

\subsection{Sources}

{\small

\begin{itemize}
\item[E1] Y. Solomon and B. Weiss, \emph{Dense forests and Danzer sets}. Primary source checked 1 October 2026. \url{https://arxiv.org/abs/1406.3807}.


\item[\hypertarget{source:1386:1}{S1}] ULAM.AI, \emph{UnsolvedMath}, supplied \texttt{problems.json}, record 1386 (GEOM-027), statement and background. The embedded literature-review date is 2026-08-17; that is the archive’s review date, not a fresh certification. Dataset landing page: \url{https://huggingface.co/datasets/ulamai/UnsolvedMath}.

\item[\hypertarget{source:1386:2}{S2}] Danzer set overview. \url{https://en.wikipedia.org/wiki/Danzer_set}. Bibliographic information imported from the supplied record.

\end{itemize}

}

\label{end:1386}
\end{document}
